\documentclass[12pt, a4paper]{article}

% 引入必要的宏包
\usepackage[UTF8]{ctex}     % 中文支持
\usepackage{amsmath, amssymb, amsthm} % 数学符号和环境
\usepackage{geometry}       % 页面设置
\usepackage{enumitem}       % 列表定制
\usepackage{graphicx}       % 图片支持
\usepackage{fancyhdr}       % 页眉页脚

% 页面边距设置
\geometry{left=2.5cm, right=2.5cm, top=2.5cm, bottom=2.5cm}

% 宏定义：方便排版选择题选项
% 使用方法：\choices{选项A}{选项B}{选项C}{选项D}
\newcommand{\choices}[4]{
    \begin{enumerate}[label=\Alph*., itemsep=0pt, parsep=0pt, topsep=5pt, labelsep=0.5em]
        \item #1
        \item #2
        \item #3
        \item #4
    \end{enumerate}
}
% 横向排列的选项 (2列)
\newcommand{\choicesTwo}[4]{
    \begin{enumerate}[label=\Alph*., itemsep=0pt, parsep=0pt, topsep=5pt, labelsep=0.5em]
        \begin{minipage}{0.45\linewidth} \item #1 \end{minipage}
        \begin{minipage}{0.45\linewidth} \item #2 \end{minipage}
        \begin{minipage}{0.45\linewidth} \item #3 \end{minipage}
        \begin{minipage}{0.45\linewidth} \item #4 \end{minipage}
    \end{enumerate}
}
% 横向排列的选项 (4列)
\newcommand{\choicesFour}[4]{
    \begin{enumerate}[label=\Alph*., itemsep=0pt, parsep=0pt, topsep=5pt, labelsep=0.5em]
        \begin{minipage}{0.22\linewidth} \item #1 \end{minipage}
        \begin{minipage}{0.22\linewidth} \item #2 \end{minipage}
        \begin{minipage}{0.22\linewidth} \item #3 \end{minipage}
        \begin{minipage}{0.22\linewidth} \item #4 \end{minipage}
    \end{enumerate}
}

\title{\textbf{《高等数学》必刷习题——基础版}}
\author{}
\date{}

\begin{document}

\section*{第十一章 二重积分（基础）}

\begin{enumerate}
    \item 设区域 $ D = \{(x, y) \mid 1 \leq x^{2} + y^{2} \leq 2\} $ ，则二重积分 $ \iint_{D} dx dy = $ \underline{\hspace{1cm}}。
    \choicesFour{$ \pi $}{$ 2\pi $}{$ 3\pi $}{$ 4\pi $}

    \item 已知 $ F(x, y) = \ln(1 + x^{2} + y^{2}) + \iint_{D} \, dx \, dy $ , 其中 $D$ 为 $xOy$ 坐标平面上的有界闭区域且 $ f(x, y) $ 在 $D$ 上连续，则 $ F(x, y) $ 在点 $(1, 2)$ 处的全微分为 \underline{\hspace{1cm}}。
    \choicesTwo
    {$ \frac{1}{3}dx+\frac{2}{3}dy $}
    {$ \frac{1}{3}dx+\frac{2}{3}dy+f(1,2) $}
    {$ \frac{2}{3}dx+\frac{1}{3}dy $}
    {$ \frac{2}{3}dx+\frac{1}{3}dy+f(1,2) $}

    \item 如果闭区域 $D$ 由 $x$ 轴，$y$ 轴及 $ x+y=1 $ 围成。则 $ \iint_{D}(x+y)^{2}d\sigma $ \underline{\hspace{1cm}} $ \iint_{D}(x+y)^{3}d\sigma $。
    \choicesFour{大于}{小于}{等于}{不确定}

    \item 设平面区域 $ D=\{(x,y)|x^{2}+y^{2}\leq1\} $ ， $ D_{1}=\{(x,y)|x^{2}+y^{2}\leq1,x\geq0,y\geq0\} $ 则下列等式一定成立的是 \underline{\hspace{1cm}}。
    \choicesTwo
    {$ \iint_{D}\sqrt{x^{2}+y^{2}}dxdy=4\iint_{D_{1}}\sqrt{x^{2}+y^{2}}dxdy $}
    {$ \iint_{D} xydxdy = 4 \iint_{D_{1}} xydxdy $}
    {$ \iint_{D} f(x,y)dxdy = 4 \iint_{D_{1}} f(x,y)dxdy $}
    {$ \iint_{D} xdxdy = 4 \iint_{D_{1}} xdxdy $}

    \item 积分区域 $D$ 为 $ x^{2}+y^{2}\leq2 $ , 则 $ \iint_{D}xd\sigma= $ \underline{\hspace{1cm}}。
    \choicesFour{$ 2\pi $}{$ \pi $}{1}{0}

    \item $D$ 是以原点为圆心，半径为 5 的圆形区域。则 $ \iint_{D} d\sigma $ \underline{\hspace{2cm}}。

    \item 设区域 $D$ 是由直线 $y = 1, x = 2$ 及 $y = x$ 所围成的三角形闭区域，则二重积分 $ \iint_{D} xy \, dx \, dy = $ \underline{\hspace{2cm}}。

    \item 区域 $D$ 是 $ x^{2} + y^{2} \leq a^{2} $ ，计算 $ \iint_{D} e^{-x^{2}-y^{2}} \, dx \, dy = $ \underline{\hspace{2cm}}。

    \item $ D = \{(x, y) \mid 1 \leq x^{2} + y^{2} \leq 4\} $ 求 $ \iint_{D} \ln(x^{2} + y^{2}) \, dx \, dy = $ \underline{\hspace{2cm}}。

    \item 设 $D$ 是矩形： $ 0 \leq x \leq a, 0 \leq y \leq b $ ，则 $ \iint_{D} dxdy = $ \underline{\hspace{2cm}}。

    \item 求二重积分 $ \iint_{D}(2x+y)dxdy $ , 其中 $D$ 是由直线 $y=3, y=x$ 及曲线 $xy=1$ 围成的闭区域。

    \item 计算二重积分 $ \iint_{D} xy \, dx \, dy $ , 其中 $D$ 是由直线 $ y = x, y = 5x $ 及 $ y = -x + 6 $ 所围成的闭区域。

    \item 计算二重积分 $ \iint_{D} xydxdy $ , 其中 $D$ 为直线 $ y = 1, x = 2 $ 及 $ y = x $ 所围成的闭区域。

    \item 求二重积分 $ \iint_{D} x^{2} dx dy $ , 其中 $D$ 是由 $ x = 0, y = 2x, y = 2 $ 所围成的区域。

    \item 求 $ I = \iint_{D} (3x + 2y) d\sigma $ ，其中 $D$ 是由两坐标轴及直线 $ x + y = 2 $ 所围成的闭区域。
\end{enumerate}


\end{document}